\documentclass[UTF8,zihao=-4]{ctexart}
\usepackage[a4paper,margin=2.1cm]{geometry}
\usepackage{array}
\usepackage{longtable}
\usepackage{hyperref}
\usepackage{verbatim}
\usepackage{graphicx}
\usepackage{tikz}
\usetikzlibrary{arrows.meta,positioning,fit,calc,shapes.geometric}
\hypersetup{hidelinks}
\setlength{\parindent}{0pt}
\setlength{\parskip}{0.55em}
\renewcommand{\arraystretch}{1.22}
\setlength{\tabcolsep}{3pt}
\emergencystretch=4em
\sloppy
\title{04 详细设计与模块化复习讲义}
\author{}
\date{}
\begin{document}
\maketitle
\tableofcontents
\newpage


整理范围：

\begin{itemize}
\item \texttt{PPT/04-01-详细设计.pdf}
\item \texttt{PPT/04-02-面向对象的模块化与信息隐藏.pdf}
\end{itemize}

交叉参考：

\begin{itemize}
\item \texttt{Exam/2012.pdf}：详细设计类图或顺序图、信息隐藏、DIP、Singleton。
\item \texttt{Exam/2013A.pdf}：职责分配、继承设计、耦合类型。
\item \texttt{Exam/2013B.pdf}：借阅者设计、LSP、内容耦合。
\item \texttt{Exam/2014.pdf}：内聚类型、职责分配、设计顺序图。
\item \texttt{Exam/2016.pdf}：支付详细设计、印记耦合、设计模式承接。
\item \texttt{Exam/2020.pdf}：库存订货规则变更、validate 方法耦合。
\item \texttt{Exam/2022.pdf}：充值优惠券与支付方式设计、包循环依赖、switch 反模式。
\item \texttt{Exam/2025.pdf}：三层顺序图、集中式/分散式/委托式控制策略、耦合与内聚评价。
\item \texttt{Review/summary-重点.pdf}、\texttt{Review/Final\_Reference.pdf}：GRASP、信息隐藏、LSP、DIP、OCP、SOFA。
\end{itemize}

说明：

\begin{itemize}
\item 本部分考查“责任放在哪里”。写答案时先说明问题，再对应原则，再给改进结构。
\item 详细设计题经常要求把需求模型转成类、接口、方法、顺序图和控制策略。
\end{itemize}

\section{二、2 个课件的主线}

{
\footnotesize
\begin{longtable}{|>{\raggedright\arraybackslash}p{0.30\textwidth}|>{\raggedright\arraybackslash}p{0.30\textwidth}|>{\raggedright\arraybackslash}p{0.30\textwidth}|}
\hline
\textbf{课件} & \textbf{主线} & \textbf{必会关键词} \\
\hline
\endfirsthead
\hline
\textbf{课件} & \textbf{主线} & \textbf{必会关键词} \\
\hline
\endhead
\texttt{04-01} 详细设计 & 从需求和架构进入类、责任和协作 & 责任、协作、GRASP、控制器、委托、Mock、集成测试 \\
\hline
\texttt{04-02} 面向对象的模块化与信息隐藏 & 用耦合、内聚和信息隐藏评价设计 & 访问耦合、继承耦合、LoD、ISP、LSP、组合、封装、OCP、DIP \\
\hline
\end{longtable}
}


\bigskip\hrule\bigskip

\section{4. 详细设计、模块化与信息隐藏}

\section{4.1 详细设计的范围}

详细设计包含：

\begin{itemize}
\item 中层设计：过程、调用关系、类协作。
\item 低层设计：数据结构、算法、类型、语句和控制结构。
\end{itemize}

输入：

\begin{itemize}
\item 用例
\item 概念模型
\item 顺序图
\item 状态图
\item 非功能需求
\item 体系结构模块规格
\item 导入和导出接口
\item 组件接口
\end{itemize}

输出：

\begin{itemize}
\item 静态结构：类、接口、关系。
\item 动态行为：对象协作和消息。
\item 状态行为：对象生命周期。
\item 实现注解：约束、异常、事务。
\item 设计理由：关键选择的依据。
\end{itemize}

\section{4.2 责任与协作}

面向对象详细设计的核心是：

\begin{itemize}
\item 责任
\item 协作
\end{itemize}

责任：

\begin{itemize}
\item 操作责任：对象要执行的任务。
\item 数据责任：对象要维护的信息。
\end{itemize}

协作：

\begin{itemize}
\item 对象之间通过消息调用完成更大的任务。
\item 一个对象不直接承担全部工作，而是把子任务委托给合适对象。
\end{itemize}

委托的意义：

\begin{itemize}
\item 降低耦合。
\item 分离关注点。
\item 让职责落在拥有信息的对象上。
\end{itemize}

\section{4.3 详细设计过程}

可按以下顺序答题：

\begin{enumerate}
\item 根据用例和系统事件识别控制入口。
\item 根据概念类图分配数据责任。
\item 根据顺序图分配操作责任。
\item 明确类、接口、属性、方法和关系。
\item 用 GRASP 原则检查职责分配。
\item 用模块化、信息隐藏和设计模式改进结构。
\item 设计单元测试和集成测试。
\end{enumerate}

\section{4.4 GRASP 原则}

\subsection{Low Coupling 低耦合}

类之间依赖越少，修改传播越小。

做法：

\begin{itemize}
\item 依赖接口。
\item 避免全局变量。
\item 避免对象链式访问。
\item 避免跨层调用。
\end{itemize}

\subsection{High Cohesion 高内聚}

一个类的职责要集中，成员之间围绕同一目标。

坏例子：

\begin{itemize}
\item 一个 \texttt{Manager} 同时处理用户、订单、支付、日志、报表。
\end{itemize}

好例子：

\begin{itemize}
\item \texttt{RechargeService} 处理充值业务。
\item \texttt{PaymentGateway} 处理支付平台交互。
\item \texttt{RechargeRecordRepository} 处理充值记录持久化。
\end{itemize}

\subsection{Information Expert 信息专家}

把责任分配给拥有完成该责任所需信息的类。

示例：

\begin{itemize}
\item 判断校园卡余额是否足够消费，责任应放在 \texttt{CampusCard} 或管理校园卡余额的服务中，而不是界面对象中。
\end{itemize}

\subsection{Creator 创建者}

B 创建 A 的条件：

\begin{itemize}
\item B 聚合或包含 A。
\item B 记录 A。
\item B 紧密使用 A。
\item B 拥有初始化 A 的数据。
\end{itemize}

示例：

\begin{itemize}
\item \texttt{RechargeService} 拥有充值金额、校园卡、支付结果等信息，负责创建 \texttt{RechargeRecord}。
\end{itemize}

\subsection{Controller 控制器}

控制器负责接收系统事件并协调业务对象。

控制器不是所有业务逻辑的堆放点。Controller 接收请求并调用 Service；Service 承担业务流程；领域对象承担自身规则。

\section{4.5 控制风格}

\subsection{集中式控制}

少数控制对象掌握大量逻辑。

优点：

\begin{itemize}
\item 决策位置集中。
\item 流程容易追踪。
\end{itemize}

缺点：

\begin{itemize}
\item 控制器臃肿。
\item 其他对象退化为数据容器。
\end{itemize}

\subsection{委托式控制}

控制器保留主流程，把子任务委托给专业对象。

这是实际项目中较稳妥的风格。

\subsection{分散式控制}

控制逻辑分散在大量小对象中。

问题：

\begin{itemize}
\item 控制流难追踪。
\item 调试和维护成本高。
\end{itemize}

\section{4.6 Parnas 信息隐藏}

Parnas 的核心思想：

\begin{quote}
模块隐藏一个重要的设计决策。
\end{quote}

模块不只是子程序。模块边界要围绕“未来容易变化的设计决策”设置。

例子：

\begin{itemize}
\item 支付平台接口会变化，应隐藏在 \texttt{PaymentGateway} 接口背后。
\item 数据库表结构会变化，应隐藏在 Repository 背后。
\item 余额更新事务策略会变化，应隐藏在 Service 内部。
\end{itemize}

\section{4.7 模块化}

模块是可以被独立设计、实现和测试的自包含单元。

面向对象中的粒度：

\begin{itemize}
\item 方法
\item 类
\item 包
\end{itemize}

模块化目标：

\begin{itemize}
\item 降低耦合。
\item 提高内聚。
\item 隐藏变化。
\item 支持独立测试。
\item 支持并行开发。
\end{itemize}

\section{4.8 耦合}

耦合表示模块之间依赖强度。耦合越低，修改传播越小。

面向对象中两类主要耦合：

\begin{itemize}
\item 访问耦合
\item 继承耦合
\end{itemize}

\subsection{访问耦合}

从高到低：

\begin{enumerate}
\item 隐式访问耦合：全局变量、单例共享状态。
\item 实现访问耦合：通过反射等方式访问内部实现。
\item 成员变量访问耦合：直接访问 public 字段。
\item 参数变量访问耦合：通过参数传递具体类型。
\item 无访问耦合。
\end{enumerate}

高风险层级：

\begin{itemize}
\item 隐式访问耦合
\item 实现访问耦合
\item 成员变量访问耦合
\end{itemize}

可接受结构：

\begin{itemize}
\item 通过参数和接口协作。
\item 不访问对方内部结构。
\end{itemize}

\subsection{Law of Demeter}

“只和直接朋友交谈”。

一个方法只调用：

\begin{itemize}
\item 自身对象的方法。
\item 参数对象的方法。
\item 本方法创建对象的方法。
\item 本对象直接字段的方法。
\end{itemize}

避免：

\begin{quote}
\footnotesize
\begin{verbatim}
order.getCustomer().getAddress().getCity()
\end{verbatim}
\end{quote}

改进方向：

\begin{quote}
\footnotesize
\begin{verbatim}
order.getCustomerCity()
\end{verbatim}
\end{quote}

或把责任放到合适对象中。

\section{4.9 接口隔离与面向接口编程}

面向接口编程：

\begin{itemize}
\item 客户端依赖抽象接口。
\item 实现类可以替换。
\item 测试时可注入 Mock。
\end{itemize}

接口隔离原则：

\begin{quote}
客户端不应被迫依赖它不使用的方法。
\end{quote}

坏接口：

\begin{quote}
\footnotesize
\begin{verbatim}
public interface CampusCardService {
    void recharge();
    void reportLost();
    void generateFinanceReport();
    void repairDevice();
}
\end{verbatim}
\end{quote}

改进：

\begin{quote}
\footnotesize
\begin{verbatim}
public interface RechargeService {
    void recharge();
}

public interface LossReportService {
    void reportLost();
}
\end{verbatim}
\end{quote}

\section{4.10 继承耦合}

继承是白盒复用，子类依赖父类内部设计。

继承耦合类型从高到低：

\begin{itemize}
\item Modification：子类修改父类语义，风险最高。
\item Refinement：子类细化父类行为。
\item Extension：子类增加新能力。
\item Nil：无继承耦合。
\end{itemize}

优先选择：

\begin{itemize}
\item 组合
\item 委托
\item 接口
\end{itemize}

\section{4.11 LSP}

LSP 要求：

\begin{quote}
子类型对象替换父类型对象后，程序正确性保持不变。
\end{quote}

契约要求：

\begin{itemize}
\item 子类前置条件不能更强。
\item 子类后置条件不能更弱。
\item 子类必须维护父类不变式。
\end{itemize}

例子：

\begin{itemize}
\item 如果父类 \texttt{Account.withdraw(amount)} 允许在余额足够时取款，子类不能额外要求“amount 必须大于 100”。
\end{itemize}

\section{4.12 组合优于继承}

组合通过对象协作复用行为。

优点：

\begin{itemize}
\item 运行时可替换。
\item 不暴露父类内部实现。
\item 更易遵守 LSP。
\item 更容易测试。
\end{itemize}

示例：

\begin{itemize}
\item \texttt{RechargeService} 组合 \texttt{PaymentGateway}，而不是继承某个支付实现类。
\end{itemize}

\section{4.13 封装}

封装层次：

\begin{itemize}
\item A：封装数据类型或抽象数据类型。
\item B：封装内部结构。
\item C：封装其他对象引用。
\item D：封装运行时类型信息。
\item E：封装潜在变化。
\end{itemize}

Iterator 模式就是封装集合内部结构，让客户端不用知道底层是 \texttt{List}、\texttt{Set} 还是 \texttt{Map}。

\section{4.14 设计原则 P1-P14}

{
\footnotesize
\begin{longtable}{|>{\raggedright\arraybackslash}p{0.30\textwidth}|>{\raggedright\arraybackslash}p{0.30\textwidth}|>{\raggedright\arraybackslash}p{0.30\textwidth}|}
\hline
\textbf{编号} & \textbf{原则} & \textbf{要点} \\
\hline
\endfirsthead
\hline
\textbf{编号} & \textbf{原则} & \textbf{要点} \\
\hline
\endhead
P1 & 全局变量有害 & 全局状态导致隐式依赖 \\
\hline
P2 & 显式表达 & 依赖、状态、约束要清楚 \\
\hline
P3 & DRY & 不重复知识和逻辑 \\
\hline
P4 & 面向接口编程 & 依赖抽象而非实现 \\
\hline
P5 & Law of Demeter & 只和直接朋友交谈 \\
\hline
P6 & ISP & 接口保持小而专用 \\
\hline
P7 & LSP & 子类可替换父类 \\
\hline
P8 & 组合优于继承 & 降低白盒复用风险 \\
\hline
P9 & 封装变化 & 把变化点藏到稳定接口背后 \\
\hline
P10 & 最小权限 & 只暴露必要访问权限 \\
\hline
P11 & OCP & 对扩展开放，对修改关闭 \\
\hline
P12 & DIP & 高层和低层都依赖抽象 \\
\hline
P13 & SRP & 一个模块只有一个变化理由 \\
\hline
P14 & SOFA & 短、单一、少参数、抽象层一致 \\
\hline
\end{longtable}
}


SOFA：

\begin{itemize}
\item Short
\item One thing
\item Few arguments
\item Abstraction level consistency
\end{itemize}

\section{4.15 集成测试与 Mock}

集成测试检查类之间的协作。

输入：

\begin{itemize}
\item 对象组合
\item 触发消息
\end{itemize}

输出：

\begin{itemize}
\item 返回值
\item 最终对象状态
\item 外部协作是否被调用
\end{itemize}

Mock 对象用于替代外部依赖：

\begin{itemize}
\item 数据库
\item 支付平台
\item 第三方 API
\item 消息队列
\end{itemize}

充值服务测试示例：

\begin{itemize}
\item 注入 Mock \texttt{PaymentGateway} 返回支付成功。
\item 注入 Mock \texttt{CampusCardRepository} 返回状态正常的校园卡。
\item 调用 \texttt{recharge}。
\item 断言余额更新、记录保存、返回结果正确。
\end{itemize}

\bigskip\hrule\bigskip

\section{对应考试模板}

\section{7.8 设计原则题模板}

题目给出代码后，按以下结构答：

\begin{enumerate}
\item 指出问题。
\item 对应原则。
\item 说明风险。
\item 给出改进。
\end{enumerate}

示例：

\begin{quote}
\footnotesize
\begin{verbatim}
问题：Controller 直接访问 Repository，并在 Controller 
中完成充值业务规则。
原则：违反分层职责、低耦合和 SRP。
风险：界面入口与数据库访问耦合，业务规则无法复用，单元测试困难。
改进：Controller 只依赖 RechargeService；
RechargeServiceImpl 注入 Repository 并封装充值流程。
\end{verbatim}
\end{quote}

\bigskip\hrule\bigskip

\section{高频考点展开}

\section{责任分配答题法}

详细设计题的核心是回答“谁负责”。按下面顺序写：

\begin{enumerate}
\item 系统事件由 Controller 接收。
\item Controller 不写业务规则，只把请求交给 Service。
\item Service 负责业务流程和事务边界。
\item 拥有数据的领域对象负责自身规则。
\item Repository 负责读取和保存。
\item 外部系统通过接口封装。
\end{enumerate}

示例：

\begin{quote}
\footnotesize
\begin{verbatim}
充值金额校验放在 RechargeService 或 Money 值对象中；
余额增加规则放在 CampusCard 中；
支付平台调用放在 PaymentGateway 中；
记录保存放在 RechargeRecordRepository 中。
\end{verbatim}
\end{quote}

\section{集中式、委托式、分散式控制}

{
\footnotesize
\begin{longtable}{|>{\raggedright\arraybackslash}p{0.23\textwidth}|>{\raggedright\arraybackslash}p{0.23\textwidth}|>{\raggedright\arraybackslash}p{0.23\textwidth}|>{\raggedright\arraybackslash}p{0.23\textwidth}|}
\hline
\textbf{控制风格} & \textbf{特征} & \textbf{优点} & \textbf{风险} \\
\hline
\endfirsthead
\hline
\textbf{控制风格} & \textbf{特征} & \textbf{优点} & \textbf{风险} \\
\hline
\endhead
集中式 & 一个控制对象掌握大量流程 & 决策集中，流程容易找 & 控制器臃肿，其他对象退化成数据容器 \\
\hline
委托式 & 控制器保留主流程，子任务交给专业对象 & 职责清晰，复用和测试较好 & 需要合理划分职责 \\
\hline
分散式 & 流程散落在许多对象中 & 单个对象较小 & 控制流难追踪，调试成本高 \\
\hline
\end{longtable}
}


考试问“适合哪种控制策略”，优先选择委托式，并说明它兼顾流程可见性和职责分配。

\section{耦合类型判断}

常见判断：

\begin{itemize}
\item 直接访问 public 字段：成员变量访问耦合，破坏封装。
\item 传入整个对象但只用少数字段：印记耦合，应只传需要的数据或引入专用参数对象。
\item 使用类型码加 switch：控制耦合，应改成多态或策略。
\item 调用 \texttt{a.getB().getC()}：消息链，违反 Law of Demeter。
\item 子类改写父类语义：违反 LSP。
\item 全局变量或共享单例状态：隐式访问耦合。
\end{itemize}

Review 高频排序：

{
\footnotesize
\begin{longtable}{|>{\raggedright\arraybackslash}p{0.45\textwidth}|>{\raggedright\arraybackslash}p{0.45\textwidth}|}
\hline
\textbf{方向} & \textbf{从差到好} \\
\hline
\endfirsthead
\hline
\textbf{方向} & \textbf{从差到好} \\
\hline
\endhead
结构化耦合 & 内容耦合 → 公共耦合 → 重复耦合 → 控制耦合 → 印记耦合 → 数据耦合 \\
\hline
结构化内聚 & 偶然内聚 → 逻辑内聚 → 时间内聚 → 过程内聚 → 通信内聚 → 功能内聚 → 信息内聚 \\
\hline
\end{longtable}
}


答题时先判定类型，再解释变更影响。耦合越高，某一模块变化越容易牵连其他模块；内聚越低，模块内部职责越杂，修改时越难定位影响范围。

\section{信息隐藏答题法}

信息隐藏不是简单把字段设为 private，而是隐藏变化点。

{
\footnotesize
\begin{longtable}{|>{\raggedright\arraybackslash}p{0.45\textwidth}|>{\raggedright\arraybackslash}p{0.45\textwidth}|}
\hline
\textbf{变化点} & \textbf{隐藏方式} \\
\hline
\endfirsthead
\hline
\textbf{变化点} & \textbf{隐藏方式} \\
\hline
\endhead
数据库存储方式 & Repository 接口 \\
\hline
支付平台 API & PaymentGateway 接口 \\
\hline
优惠券计算方式 & CouponPolicy 或 DiscountStrategy \\
\hline
订货规则 & ReorderPolicy \\
\hline
集合内部结构 & Iterator \\
\hline
创建具体产品 & Factory Method 或 Abstract Factory \\
\hline
\end{longtable}
}


\section{LSP 判断模板}

子类能替换父类的条件：

\begin{itemize}
\item 不加强父类前置条件。
\item 不削弱父类后置条件。
\item 维护父类不变式。
\item 不改变父类方法的承诺语义。
\end{itemize}

正方形继承长方形是经典反例：长方形允许独立设置长和宽，正方形要求长宽相等，子类加强了不变式，替换后会破坏调用方预期。

\section{自测题与答案}

\section{题 1：什么是信息专家原则？}

答案：

信息专家原则要求把责任分配给拥有完成该责任所需信息的类。这样能减少跨对象取数据，降低耦合并提高内聚。例：校园卡余额是否能更新，应由 \texttt{CampusCard} 或管理校园卡余额的业务对象判断。

\section{题 2：为什么 Controller 不应直接访问 Repository？}

答案：

Controller 属于请求入口，直接访问 Repository 会把界面入口、业务规则和数据访问耦合在一起。结果是业务逻辑难复用、测试困难、事务边界不清。改进方式是 Controller 依赖 Service，Service 依赖 Repository。

\section{题 3：\texttt{switch(type)} 分发行为违反什么原则？}

答案：

它违反 OCP 和低耦合，新增类型时要修改原有方法。改进方式是使用多态或策略模式，把不同类型行为放入不同实现类。

\section{题 4：什么是印记耦合，如何消除？}

答案：

印记耦合是调用方把一个复杂数据结构传给被调用方，而被调用方只使用其中一部分数据。消除方式是只传需要的字段，或定义专用参数对象，避免暴露完整结构。

\section{题 5：什么是 Law of Demeter？}

答案：

Law of Demeter 要求对象只和直接朋友通信。一个方法只调用自身、参数、本方法创建对象和直接字段的方法，避免长消息链。这样能降低对象之间的间接依赖。

\section{题 6：Mock 在详细设计和测试中有什么作用？}

答案：

Mock 用来替代外部依赖，使测试集中在当前对象协作上。支付、数据库、第三方 API 都可用 Mock 替代，从而验证 Service 是否正确调用依赖、处理返回结果并维护对象状态。

\bigskip\hrule\bigskip

\section{快速背诵清单}

\section{详细设计}

\begin{itemize}
\item 核心：责任与协作。
\item GRASP：低耦合、高内聚、信息专家、创建者、控制器。
\item 信息隐藏：模块隐藏重要设计决策。
\item SOFA：短、单一、少参数、抽象层一致。
\end{itemize}
\end{document}

